% Source-derived chapter 08: Proof of the main theorem on MCG-finite representations
% Original source: canonical-representations-surface-groups
\chapter{Proof of the main theorem on MCG-finite representations}
\label{canonical-representations-surface-groups-chapter-08}
\source{canonical-representations-surface-groups}

\begin{remark}[Source section]
\label{canonical-representations-surface-groups-chapter-08-source}
\source{canonical-representations-surface-groups}
\source{canonical-representations-surface-groups}
This chapter reproduces the corresponding section of the retrieved
LaTeX source, including its definitions, statements, examples, and
proofs. It has not yet been translated into Lean.
\notready
\end{remark}

% The source section title was: Proof of the main theorem on MCG-finite representations
\label{section:main-proof}
\source{canonical-representations-surface-groups}

We have finally developed the tools to verify our main result, \autoref{theorem:finite-image}: that  MCG-finite representations of rank $r<\sqrt{g+1}$ have finite image.
The proof is fairly involved, so the reader may find it useful to refer to the sketch in \autoref{subsection:outline-of-proof}. Briefly, we 
%and so we sketch the idea.
%At its core, we aim to apply \cite[Theorem 1.2.5]{LL:geometric-local-systems},
%which states that $\mathscr O_K$ local systems of rank at most $2\sqrt{g+1}$
%on an analytically very general curve,
%which underlies a 
%a complex polarizable variation of Hodge structure
%has finite monodromy.
%Starting with a unitary MCG-finite representation, we wish to verify the above
%mentioned hypotheses.
%We 
start by reviewing the notion of cohomological rigidity in
\autoref{subsection:background-rigidity}
and prove the necessary rigidity results in
\autoref{subsection:rigidity}, using our main vanishing result,
\autoref{theorem:unitary-rigid}.
We next deduce the integrality of MCG-finite representations of low rank in
\autoref{subsection:integrality}
using a result of 
Klevdal-Patrikis \cite{klevdalP:g-rigid-local-systems-are-integral}, which builds on work of Esnault-Groechenig \cite{esnault2018cohomologically} and ultimately relies on input from the Langlands program and known cases of the companion conjectures.
We then combine the above ingredients with input from non-abelian Hodge theory to deduce that low rank unitary MCG-finite representations
have finite image in \autoref{subsection:finiteness}.
We generalize this to the semisimple but non-unitary case in
\autoref{subsection:semisimple-case}, again using non-abelian Hodge theory.
Finally, in \autoref{subsection:non-semisimple} we bootstrap these results to the general, non-semisimple case,
using \autoref{theorem:asymptotic-putman-wieland} on the Putman-Wieland conjecture.


\subsection{Background on Cohomological rigidity}
\label{subsection:background-rigidity}
\source{canonical-representations-surface-groups}

We now define the notion of cohomological rigidity, which detects whether a
representation has any first order deformations.

\begin{definition}[Cohomological rigidity]
\source{canonical-representations-surface-groups}
\notready\label{definition:rigidity}
\source{canonical-representations-surface-groups}
    Let $\overline{X}$ be a smooth projective variety, $Z\subset \overline{X}$ a
    strict normal crossings divisor, and $X=\overline{X}\setminus Z$. Let $G$ be
    a reductive group over $\mathbb{C}$ and $$\rho: \pi_1(X)\to G(\mathbb{C})$$
    a homomorphism such that the monodromy at infinity is quasi-unipotent.   
%    The representation $\rho$ is said to be \emph{rigid} if for any homomorphism $$\rho': \pi_1(X)\to G(\mathbb{C}[[t]])$$ with $\rho'|_{t=0}=\rho$ and with quasi-unipotent monodromy at infinity, there exists $g\in G(\mathbb{C}[[t]])$ with $g|_{t=0}=\on{id}$ such that $$g\rho'g^{-1}=\rho\otimes \mathbb{C}[[t]].$$ That is, $\rho$ is an isolated point of the appropriate character variety.
     The representation $\rho$ is said to be \emph{cohomologically rigid} if
    $$H^1(\overline{X}, j_{!*}\on{ad}\rho)=0,$$ where $j: X\hookrightarrow
    \overline{X}$ is the natural inclusion, $j_{!*}$ is the intermediate
    extension, and $\on{ad}\rho$ is defined as in \autoref{notation:adjoint}.
    
    If $\mathbb{V}$ is the local system on $X$ associated to $\rho$, we call $\mathbb{V}$ cohomologically rigid if $\rho$ is.
\end{definition}
\begin{remark}
\source{canonical-representations-surface-groups}
\notready
	\label{remark:}
\source{canonical-representations-surface-groups}
	In the case $\rho$ is irreducible, see \cite[Proposition 4.7, Remark
    4.8]{klevdalP:g-rigid-local-systems-are-integral} for a moduli-theoretic
    interpretation of cohomological rigidity. 
    This says that $\rho$ is a smooth isolated point of the appropriate character variety, parametrizing representations with {fixed} local monodromy at infinity. Note that if $\overline{X}$ is a curve, $j_{!*}=j_*.$
\end{remark}

We will typically prove cohomological rigidity by computing 
$H^1(X, \on{ad}\rho)$, via the next lemma.
\begin{lemma}
\source{canonical-representations-surface-groups}
\notready
	\label{lemma:vanishing-h1-implies-rigidity}
\source{canonical-representations-surface-groups}
		In the setting of \autoref{definition:rigidity}, assume $H^1(X, \on{ad}\rho)
	= 0$. Then $\rho$ is
	cohomologically rigid.
\end{lemma}
\begin{proof}
	This follows from the identification 
	$H^1(\overline{X}, j_{!*}\on{ad}\rho)\simeq H^1(U, a_*\on{ad}\rho)$ for
	$U \subset \overline{X}$ a certain dense open subscheme containing $X$
	with $a: X \to U$ the inclusion \cite[Remark
	2.4]{esnault2018cohomologically} (see also \cite[Remark
	4.8]{klevdalP:g-rigid-local-systems-are-integral}).
	The Leray spectral sequence gives an injection 
	$H^1(U, a_*\on{ad}\rho) \hookrightarrow H^1(X, \on{ad}\rho)$, implying
	the claim.
\end{proof}
\begin{definition}
\source{canonical-representations-surface-groups}
\notready
If $H^1(X, \on{ad}\rho)=0$, we say that $\rho$ is \emph{strongly cohomologically rigid}.	
\end{definition}
\subsection{Cohomological rigidity and unitary MCG-finite representations}
\label{subsection:rigidity}
\source{canonical-representations-surface-groups}

Let $\rho: \pi_1(\Sigma_{g,n})\to \on{GL}_r(\mathbb{C})$ be a unitary MCG-finite
representation of rank $<\sqrt{g+1}$. We will study $\rho$ by associating to $\rho$ a certain unitary local system on a well-chosen versal family of curves. Crucially, this unitary local system will be cohomologically rigid.
This will follow from our main vanishing result
\autoref{theorem:unitary-rigid}.

Throughout this section, we use notation as in 
\autoref{notation:versal-family}. In particular, $\pi:\mathscr{C}\to
\mathscr{M}$ is a versal family of $n$-pointed curves of genus $g$,
$\pi^\circ: \mathscr{C}^\circ \to \mathscr{M}$ is the associated family of
punctured curves, $m \in \mathscr M$ is a basepoint, and $C^\circ = \mathscr
C_m^\circ$.

\begin{proposition}
\source{canonical-representations-surface-groups}
\notready
	\label{proposition:cohomologically-rigid}
\source{canonical-representations-surface-groups}
	Let $\mathbb{V}$ be a $\on{GL}_r$ (respectively, ~$\on{PGL}_r$)-local system on
$\mathscr{C}^\circ$ with $r<\sqrt{g+1}$. Suppose that for $m\in \mathscr{M}$,
$\mathbb{V}|_{{C}^\circ}$ is (respectively, is the projectivization
of) an irreducible,
unitary local system. Then $\mathbb{V}$ is strongly cohomologically rigid. In
particular, $\mathbb V$ is cohomologically rigid.
\end{proposition}
\begin{proof}
	By \autoref{lemma:vanishing-h1-implies-rigidity}, it is enough to show
	$\mathbb V$ is strongly cohomologically rigid.
	We let $\on{ad}\mathbb V$ denote the adjoint local system as in
	\autoref{notation:adjoint}.
	We will check 
	$H^1(\mathscr C^\circ, \on{ad} \mathbb V) = 0$.
	Using the Leray spectral sequence associated to the map $\pi^\circ$, it suffices to show that  
	$$H^0(\mathscr M, R^1 \pi^\circ_* \on{ad} \mathbb V) = 
	H^1(\mathscr M, \pi^\circ_* \on{ad} \mathbb V) = 0.$$
	We have 
	$H^0(\mathscr M, R^1 \pi^\circ_* \on{ad} \mathbb V) = 0$
	by \autoref{theorem:unitary-rigid} (taking $A=\mathbb{C}$), as $\rk \on{ad}\mathbb{V}=r^2-1<g$ and $\on{ad}\mathbb{V}|_{{C}^\circ}$ is unitary by assumption.

	To conclude, it is enough to show $\pi^\circ_* \on{ad} \mathbb V = 0$.
	This is follows from Schur's lemma because
	$\mathbb{V}|_{{C}^\circ}$ is (the projectivization of) an
	irreducible local system, and so $\on{ad}\mathbb{V}|_{{C}^\circ}$ has no $\pi_1({C}^\circ)$-invariants.
\end{proof}
%\begin{remark}
%	\label{remark:tangent-space}
%	For $\mathbb{V}$ as in \autoref{proposition:cohomologically-rigid},
%	corresponding to the local systems $\mathbb W$,
%	we showed in the proof of \autoref{proposition:cohomologically-rigid} that 
%	$H^1(\mathscr C^\circ, \on{ad} \mathbb{V})$
%	injects into  $H^0(\mathscr M, R^1 \pi^\circ_* \on{ad}\mathbb{V})$.
%	This crucially used irreducibility of $\mathbb{V}|_{\mathscr{C}^\circ_m}$.
%	The group 
%	$H^1(\mathscr C^\circ, \on{ad} \mathbb V)$
%	can be identified as the tangent space at $[\mathbb V]$ to the stack of
%	representations of $\pi_1(\mathscr C^\circ, c)$ up to
%	conjugacy.
%	Therefore, at least when we restrict to irreducible local systems
%	$\mathbb$,
%	\autoref{proposition:cohomologically-rigid} implies the following: 
%	suppose $\rho$ is a unitary representation 
%	whose conjugacy class has finite orbit under $\on{Mod}_{g,n}$
%	and $\dim \rho < \sqrt {g}$ then $\mathbb W$ is an isolated
%	reduced periodic point of the action of $\on{Mod}_{g,n}$ on the
%	above described stack of representations.
%
%	We would like to moreover say that $H^0(\mathscr M, R^1 \pi^\circ_* \on{ad}\mathbb W)$ 
%	describes ``deformations of representations of
%	$\pi_1(\Sigma_{g,n},x)$ with fixed determinant which are invariant under the $\pi_1(\mathscr
%	M,m)$ action up to conjugacy.''
%	This works fine when $\rho$ is irreducible, via the equivalence between
%	this an the above described algebraic stack of projective
%	representations of $\pi_1(\mathscr C^\circ,c)$.
%
%	However, when $\rho$ is reducible,
%	we cannot canonically extend it to a projective local
%	system $\mathbb W$ on $\mathscr C^\circ$ as in
%	\autoref{lemma:geometric-mcg-rep-construction}, which uses Schur's
%	lemma.
%	Indeed, for reducible $\rho$, there is no such canonical extension,
%	and there may either be
%	many such extensions or no extensions whatsoever.
%\begin{proof}[Proof sketch]
%	Let $\Gamma\subset \on{Mod}_{g,n}$ be a finite index subgroup of
%	the stabilizer of $[\mathbb{W}]$, and let $\mathscr{M}$ be the
%	corresponding cover of $\mathscr{M}_{g,n}$ and $\mathscr C^\circ$ the
%	universal punctured curve over $\mathscr M$. Then 
%	$H^1(\mathscr C^\circ, \on{ad} \mathbb W)$ is the tangent space to
%	the stack parameterizing $\pi_1(\mathscr M)$ invariant representations
%	of $\pi_1(\Sigma_{g,n})$ with fixed determinant.
%	By \autoref{theorem:unitary-rigid},
%	this tangent space is trivial, so $[\mathbb W$ is an isolated point.
%\end{proof}
%\end{remark}
\subsection{Verifying integrality}
\label{subsection:integrality}
\source{canonical-representations-surface-groups}

The next step of the proof is to use the cohomological rigidity provided by
\autoref{proposition:cohomologically-rigid} to deduce integrality of MCG-finite unitary representations.

\begin{definition}
\source{canonical-representations-surface-groups}
\notready
	\label{definition:}
\source{canonical-representations-surface-groups}
	Let $G$ be a group scheme over $\mathbb{Z}$. We say a representation $\pi_1(X, x) \to G(\mathbb C)$ is {\em integral}
if it is conjugate to a representation which factors through $G(\mathscr
O_K)$ for some number field $K$.
A local system on a connected space $X$ is integral if the corresponding monodromy representation is.
\end{definition}
We will primarily be concerned with integrality with respect to the group
schemes $G = \on{PGL}_r$ and $G = \on{GL}_r$.
The basic idea to show our local system is integral is to apply the main result of
\cite{klevdalP:g-rigid-local-systems-are-integral}, which builds on \cite{esnault2018cohomologically} and ultimately relies on Lafforgue's work on the Langlands program for function fields and consequent work on Deligne's companion conjectures.
To apply this result, we need to know the monodromy of our cohomologically rigid
local system $\mathbb{V}$ around boundary components of a strict normal crossings compactification of $\mathscr{C}^\circ$
is quasi-unipotent, which we will verify using the following result.

\begin{proposition}[\protect{~\cite[Proposition 2.4]{aramayona2016rigidity}}]
\source{canonical-representations-surface-groups}
\notready
			\label{proposition:quasi-unipotent-dehn}
\source{canonical-representations-surface-groups}
	Suppose that $g \geq 3$, $\gamma \subset \Sigma_{g,n}$ is a simple closed curve, 
	and $\Gamma\subset \on{Mod}_{g,n}$ is a 
finite index subgroup. Let $\delta_\gamma\in\on{Mod}_{g,n}$ be the Dehn twist about $\gamma$. Let $\rho: \Gamma\to \on{GL}_r(\mathbb{C})$ be a representation. Then for $m$ such that $\delta_\gamma^m\in \Gamma$, $\rho(\delta_\gamma^m)$ is quasi-unipotent.
\end{proposition}
%\begin{proposition}
%	Suppose $g\geq 2$, $\gamma \subset \Sigma_{g,n+1}$ a simple closed curve, 
%	$\Sigma_{g,n+1} - \gamma$ has some connected component of genus at least
%	$2$,
%	and $\delta_\gamma\in \widetilde{\Gamma}_\rho\subset \on{Mod}_{g,n+1}$ is a 
%power of a Dehn twist. Then $\tilde\rho(\delta_\gamma)$ is quasi-unipotent.
%\end{proposition}
%\begin{proof}
%The proof is given in the proof of \cite[Proposition 2.4]{aramayona2016rigidity}.
%Note that the statement of \cite[Proposition 2.4]{aramayona2016rigidity}
%assumes $g \geq 3$, with $\Gamma$ and simple closed curve, but the proof only
%assumes that one of the components of $\Sigma_{g,n+1} - \gamma$ has genus at
%least $2$.
%\end{proof}

We now verify that the projectivization of a unitary MCG-finite representation is integral.
\begin{lemma}
\source{canonical-representations-surface-groups}
\notready
	\label{lemma:defined-over-number-field}
\source{canonical-representations-surface-groups}
	Let $g\geq 3$ be an integer, and let $\rho: \pi_1(\Sigma_{g,n})\to \on{GL}_r(\mathbb{C})$ be an irreducible, unitary, MCG-finite representation, with $r<\sqrt{g+1}$. Then the composition $$\mathbb{P}\rho: \pi_1(\Sigma_{g,n})\overset{\rho}{\to} \on{GL}_r(\mathbb{C})\to \on{PGL}_r(\mathbb{C})$$ is integral.
\end{lemma}
\begin{proof}
By \autoref{lemma:geometric-mcg-rep-construction}, there exists a scheme $\mathscr{M}$ with a finite \'etale map $\mathscr{M}\to \mathscr{M}_{g,n}$, with associated family of punctured curves $\pi^\circ: \mathscr{C}^\circ\to \mathscr{M}$, and a representation $$\widetilde{\rho}: \pi_1(\mathscr{C}^\circ)\to \on{PGL}_r(\mathbb{C})$$ agreeing with $\mathbb{P}\rho$ on restriction to a fiber of $\pi^\circ$. It suffices to show that $\widetilde{\rho}$ is integral.  
	By \autoref{proposition:cohomologically-rigid},
	$\widetilde{\rho}$ is cohomologically rigid.
	Because the abelianization of $\on{PGL}_r(\mathbb C)$ is
	trivial, it follows from
\cite[Theorem 1.2]{klevdalP:g-rigid-local-systems-are-integral}	that $\widetilde\rho$
is integral once we verify that $\widetilde\rho$ has quasi-unipotent local monodromy around the boundary components of
	a good (i.e.~strict normal crossing) compactification of $\mathscr{C}^\circ$.

One may construct a strict normal crossing compactification as follows.
$\mathscr{C}^\circ$ is a finite \'etale cover of $\mathscr{M}_{g,n+1}$ by
construction; let $\overline{\mathscr{M}_{g, n+1}}'$ be a strict normal crossing
compactification of $\mathscr{M}_{g,n+1}$ obtained by blowing up boundary strata
of the Deligne-Mumford compactification $\overline{\mathscr{M}_{g,n+1}}$ (which
is only a normal crossing compactification, and is not in general strict). Let
$\overline{\mathscr{C}}$ be the normalization of $\overline{\mathscr{M}_{g,
n+1}}'$ in the function field of $\mathscr{C}^\circ$. By
\autoref{proposition:quasi-unipotent-dehn}, it suffices to check that the local
monodromy about the boundary components of $\overline{\mathscr{C}}$ correspond
to products of commuting Dehn twists about simple closed curves, under the
identification of $\pi_1(\mathscr{C}^\circ)$ with a subgroup of
$\on{PMod}_{g,n+1}=\pi_1(\mathscr{M}_{g,n+1})$. 
This is because a product of
commuting quasi-unipotent matrices is quasi-unipotent. To check the claim about
local monodromy about boundary components, it suffices to check the
corresponding claim for $\overline{\mathscr{M}_{g, n+1}}'$. But this is
\cite[Lemma 2.1.1]{li2020surface}.
\end{proof}

The next result shows we can lift integrality from the projectivization of a unitary irreducible MCG-finite representation to the original representation. 

\begin{lemma}
\source{canonical-representations-surface-groups}
\notready
	\label{lemma:gl-defined-over-number-field}
\source{canonical-representations-surface-groups}
	Let $g\geq 3$ be an integer, and let $\rho: \pi_1(\Sigma_{g,n})\to \on{GL}_r(\mathbb{C})$ be an irreducible, unitary, MCG-finite representation, with $r<\sqrt{g+1}$. Then $\rho$ is integral.
	\end{lemma}
\begin{proof}
By \autoref{proposition:1-dim-finite-image}, $\det\rho$ has finite image, and
hence $\rho$ factors through $G(\mathbb{C})$, where $G\subset \on{GL}_r$ is a
flat affine $\mathbb{Z}$-group scheme containing $\on{SL}_r$ with finite index (for
example, we can take $G=\det^{-1}(\mu_m)$, where $\det\rho$ factors through
$\mu_m$). Note that the natural morphism $G\to \on{PGL}_r$ is finite. After
conjugating $\rho$ by some matrix we may assume $\mathbb{P}\rho$ factors through
$\on{PGL}_r(\mathscr{O}_K)$ for some number field $K$, by \autoref{lemma:defined-over-number-field}.

As $\pi_1(\Sigma_{g,n})$ is finitely generated, it suffices to show that for a generator $\alpha\in \pi_1(\Sigma_{g,n})$, there exists some finite extension $K'/K$ such that the matrix $\rho(\alpha)$ lies in $G(\mathscr{O}_{K'})$. 
The obstruction to lifting any given $\mathscr O_K$ point of $\on{PGL}_r$ to
$G$ is a $\ker(G\to \on{PGL}_r)$ torsor, which becomes trivial on some finite
flat extension, and hence over $\mathscr O_{K'}$ for some finite $K'/K$.
%%Here is an alternate, longer proof:
%	The second statement follows from the first, because we verified
%	$\mathbb P$ is integral in \autoref{lemma:defined-over-number-field}, 
%	and $\mathbb W$ has finite determinant by the final statement of
%	\autoref{proposition:gl-lift} that
%	there is a group $G$ containing
%	$\on{SL}_r(\mathbb C)$ with finite index so that $\rho'$ factors through
%	$G$.
%
%	Observe that the determinant of $G$ is contained in a finite group
%	$\mu_m(\mathbb C)$, and so $H := \on{det}^{-1}(\mu_m) \subset \on{GL}_r$
%	defines an algebraic group with $G \subset H(\mathbb C)$.
%	Observe that $\mu_{mr} \simeq \ker(H \to \on{PGL}_r)$
%	After possibly enlarging $K$, we may assume all $mr$th roots of unity
%	lie in $K$.
%
%	Viewing $\mathbb W$ as a representation $\widetilde{\rho} : \pi_1(\mathscr C^\circ,
%	c) \to \on{PGL}_r(\mathscr O_K)
%	\subset \on{PGL}_r(K) \subset \on{PGL}_r(\mathbb C)$, we claim there
%	exists a finite extension $K'$ of $K$ so that
%	$\mathbb V$, viewed as a representation $\rho': \pi_1(\mathscr C^\circ,
%	c) \to H(\mathbb C)$
%	factors through $H(K')$.
%	Indeed, since we are assuming all $mr$th roots of unity lie in $K$,
%	every element of $\on{PGL}_r(K)$ lifts to an element of $H(K)$, as the
%	obstruction lives in $H^1(K, \mu_{mr}) \simeq K^\times/(K^\times)^{mr}$.
%	Therefore, the image of any generator of $\pi_1(\mathscr C^\circ)$ in
%	$\on{PGL}_r(K)$ lifts to an element of $H(K)$ after adjoining a
%	particular $mr$th root of a particular element to $K$.
%	Note that $\pi_1(\mathscr C^\circ)$ is finitely generated since it
%	has finite index in $\on{Mod}_{g,n+1}
%	\simeq \pi_1(\mathscr M_{g,n+1})$, which is finitely generated \cite[Theorem 5.7]{farbM:a-primer}.
%	Hence, we now take our desired $K'$ to be the finite field extension of
%	$K$ obtained by adjoining finitely many such $n$th roots of $K$ so that
%	each generator of $\pi_1(\mathscr C^\circ)$ maps to an element of
%	$\on{PGL}_r(K)$ which lifts to some element of $H(K')$.
%	Since each generator maps to an element of
%	$\on{PGL}_r(K)$ which has some lift in $H(K')$ and the set of possible
%	lifts is a torsor for $\mu_r(\mathbb C)$, which we are assuming is
%	contained in $K\subset K'$, every such lift lies in $H(K')$.
%	Therefore, $\rho'$ indeed factors through $H(K')$.
%
%	To conclude, we wish to show that $\rho'$ moreover factors through
%	$H(\mathscr O_{K'})$.
%	By assumption, $\widetilde{\rho}$ factors through $\on{PGL}_r(\mathscr
%	O_{K}) \subset \on{PGL}_r(\mathscr
%	O_{K'})$.
%	Since $H \to \on{PGL}_r$ is a finite map, the valuative criterion for
%	properness implies that for each $g \in \pi_1(\mathscr C^\circ, c)$,
%	$\rho'(g) \in H(K')$ in fact lies in $H(\mathscr O_{K'})$, as its image
%	in $\on{PGL}_r(K')$ factors through $\on{PGL}_r(\mathscr O_{K'})$.
%	This uniqueness then implies the resulting map
%	$\pi_1(\mathscr C^\circ, c) \to H(\mathscr O_{K'})$ is a group
%	homomorphism, this produces the desired integral structure on $\mathbb V$.
\end{proof}

\subsection{Verifying finiteness}
\label{subsection:finiteness}
\source{canonical-representations-surface-groups}
We now show that a unitary MCG-finite representation $\rho$ of low rank has finite image. We have already shown that, in the irreducible setting, such representations are defined over $\mathscr{O}_K$. We will use the techniques of \cite{LL:geometric-local-systems} to deduce that they have finite monodromy. To do so, we will argue that for each embedding $\tau: \mathscr{O}_K\hookrightarrow\mathbb{C}$, $\rho\otimes_{\mathscr{O}_K, \tau}\mathbb{C}$ has unitary monodromy. We know this for one such $\tau$ but not for the rest.  
To prove unitarity for all such $\tau$, we will use \cite[Theorem
1.2.12]{LL:geometric-local-systems}, and to verify its hypotheses we will need
to use non-abelian Hodge theory to construct certain complex PVHS.

\begin{proposition}
\source{canonical-representations-surface-groups}
\notready\label{proposition:unitary-implies-finite}
\source{canonical-representations-surface-groups}
Let  $$\rho:\pi_1(\Sigma_{g,n})\to
\on{GL}_r(\mathbb{C})$$ be a unitary MCG-finite representation with $r<\sqrt{g+1}$. Then $\rho$ has finite image.
\end{proposition}
\begin{proof}
As $\rho$ is unitary, it is semisimple. As an irreducible sub-representation of an MCG-finite representation is MCG-finite by \autoref{proposition:basic-properties}, we may without loss of generality assume $\rho$ is irreducible.

If $g< 3$, $r\leq 1$, and so we may conclude by \autoref{proposition:1-dim-finite-image}. Thus we may assume $g\geq 3$. By \autoref{lemma:gl-defined-over-number-field}, we may assume $\rho$ factors through $\on{GL}_r(\mathscr{O}_K)$ for some number field $K$ and some embedding $\iota: \mathscr{O}_K\hookrightarrow \mathbb{C}$.
We now adjust notation so that $\rho$ denotes this $\mathscr O_K$-representation and $\rho_\iota=\rho\otimes_{\mathscr{O}_K, \iota}\mathbb{C}$
denotes our original MCG-finite representation. By \autoref{corollary:converse-to-MCG-rep}, there exists a punctured versal family $\mathscr C^\circ \to \mathscr M$
and a representation $\rho'_\iota: \pi_1(\mathscr C^\circ) \to
\on{GL}_r(\mathbb{C})$ with finite determinant, restricting to $\rho_\iota$ on
the fundamental group of a fiber of $\pi^\circ$, i.e.~ for $C^\circ$ a fiber of
$\pi^\circ$, $\rho'_\iota|_{\pi_1(
C^\circ)}=\rho_\iota$. Note that $\rho'_\iota$ is irreducible (as $\rho_\iota$
is irreducible), and strongly cohomologically rigid, by \autoref{proposition:cohomologically-rigid}.

Choose any embedding $\tau: \mathscr O_K \hookrightarrow \mathbb C$ and let
$\rho_\tau := \rho\otimes_{\mathscr{O}_K, \tau}\mathbb{C}$. Note that $\rho_\tau$ is MCG-finite, as the same is true for $\rho_\iota$. 
We aim to show $\rho_\tau$ is unitary.

Choose $\sigma: \mathbb{C}\overset{\sim}{\to} \mathbb{C}$ so that $\sigma\circ\iota=\tau$, and set $\rho_\tau'=\rho_{\iota}'\otimes_{\mathbb{C}, \sigma} \mathbb{C}$. Note that $\rho_\tau'|_{\pi_1(
C^\circ)}=\rho_\tau$.
Now observe that $\rho_\tau'$ is strongly cohomologically rigid, as the same is true for $\rho_\iota'$, and strong cohomological rigidity (being a cohomological condition) is preserved by automorphisms of $\mathbb{C}$. Moreover $\rho_\tau'$ is irreducible with finite determinant, as the same is true for $\rho_\iota'$.

We next show $\rho_\tau$ is unitary.
By \cite[Theorem 1.2.12]{LL:geometric-local-systems}, 
in order to show $\rho_\tau$ is unitary, it is enough to show $\rho'_\tau$ underlies
a complex PVHS on $\mathscr{C}^\circ$.
Since $\rho'_\tau$ is strongly cohomologically rigid and irreducible with finite determinant,
this follows from \autoref{lemma:rigid-reps}.

We are thus in the following situation: $\rho: \pi_1(\Sigma_{g,n})\to \on{GL}_r(\mathscr{O}_K)$ is a representation such that for each $\tau: \mathscr{O}_K\hookrightarrow \mathbb{C}$, $\rho\otimes_{\mathscr{O}_K, \tau} \mathbb{C}$ is unitary. Such a representation has finite image by \cite[Lemma 7.2.1]{LL:geometric-local-systems}.
\end{proof}

\subsection{Reduction from the semisimple case to the unitary case}
\label{subsection:semisimple-case}
\source{canonical-representations-surface-groups}
In this section we will prove \autoref{theorem:finite-image} for semisimple
representations.
The idea is to first deform our representation to a complex PVHS via
non-abelian Hodge theory. Then we will argue that this complex PVHS is in fact unitary, using the results of \cite{LL:geometric-local-systems}, so that we
can apply \autoref{proposition:unitary-implies-finite}. Finally we will use the rigidity properties of unitary MCG-finite representations to argue 
that our deformation was in fact trivial---the deformed representation was in fact the one we started with.

\begin{lemma}
\source{canonical-representations-surface-groups}
\notready\label{lemma:infinitesmal-rigidity}
\source{canonical-representations-surface-groups}
Let $\pi^\circ: \mathscr{C}^\circ\to \mathscr{M}$ be a punctured versal family
of curves of genus $g$. Let $c\in \mathscr{C}^\circ$ be a point, set
$m=\pi^\circ(c)$, and let $C^\circ=\mathscr{C}^\circ_m$. Let $$\rho_\infty: \pi_1(\mathscr{C}^\circ, c)\to GL_r(\mathbb{C}[[t]])$$ be a representation with $r<\sqrt{g+1}$ and constant determinant, and let $\rho_0=\rho_\infty\otimes\mathbb{C}$. Suppose $\rho_0|_{\pi_1(C^\circ, c)}$ is unitary. Then $\rho_\infty|_{\pi_1(C^\circ, c)}$ is conjugate to a constant representation.
\end{lemma}
\begin{proof}
	We apply \autoref{prop:stacky-split-constancy} to the short exact sequence $$1\to \pi_1(C^\circ, c)\to \pi_1(\mathscr{C}^\circ, c)\to \pi_1(\mathscr{M},m)\to 1.$$ We must verify that for any $n$ and any deformation $\gamma: \pi_1(\mathscr{C}^\circ, c)\to GL_r(\mathbb{C}[t]/t^{n+1})$ of $\rho_0$ with $\gamma|_{\pi_1(C^\circ, c)}$ constant, we have $$H^0(\mathscr{M}, R^1\pi_*\on{ad}(\gamma))=0.$$ But this follows from \autoref{theorem:unitary-rigid}. Indeed, $$\rk_{\mathbb{C}[t]/t^{n+1}} \on{ad}(\gamma)=r^2-1<g$$ as $r<\sqrt{g+1}$ by assumption.
\end{proof}


\begin{lemma}
\source{canonical-representations-surface-groups}
\notready\label{lemma:global-rigidity}
\source{canonical-representations-surface-groups}
Let $\pi^\circ: \mathscr{C}^\circ\to \mathscr{M}$ be a punctured versal family
of curves of genus $g$. Let $c\in \mathscr{C}^\circ$ be a point, set
$m=\pi^\circ(c)$, and let $C^\circ=\mathscr{C}^\circ_m$. Let $X$ be a
finite-type connected $\mathbb{C}$-scheme, and let $$\rho:
\pi_1(\mathscr{C}^\circ, c)\to \on{GL}_r(\mathscr{O}_X(X))$$ be a representation
with $r<\sqrt{g+1}$ and constant determinant. For a closed point $x\in X$ with residue field
$\kappa(x)=\mathbb{C}$, let $\rho_x:=\rho\otimes_{\mathscr{O}_X(X)}\kappa(x)$.
Suppose there exists a closed point $y\in X$ such that $\rho_y|_{\pi_1(C^\circ,
c)}$ is unitary. Then for each
closed point $x\in X$, $\rho_{x}|_{\pi_1(C^\circ, c)}$ is conjugate to $\rho_y|_{\pi_1(C^\circ, c)}$.
\end{lemma}
\begin{proof}
Without loss of generality $X$ is a reduced (not necessarily irreducible) connected curve, as any closed point $x\in X$ is connected to $y$ by such a curve. We claim it suffices to prove the statement if the parameter space $X$ is smooth and irreducible. 

To reduce to this case, suppose $X$ is neither smooth nor irreducible, and let $\tilde Y$ be the normalization of the component $Y$ of $X$ containing $y$, and $\iota:
\tilde Y\to X$ the natural map. By the case of smooth irreducible curves, $\iota^*\rho|_{\pi_1(C^\circ, c)}\otimes \kappa(y')$ is independent of $y'\in \tilde Y(\mathbb{C})$, and thus $\rho|_{\pi_1(C^\circ, c)}\otimes \kappa(y'')$ is independent of $y''\in Y$ as well, and in particular unitary for all $y''$. Let $Z$ be a connected component of the closure of $X-\on{im}(\iota)$. As $X$ is connected, $Z$ intersects $Y$ non-trivially, say at some point $z\in Y$. But as $z$ is in $Y$, $\rho_z|_{\pi_1(C^\circ, c)}$ is unitary, and hence we are done by induction on the number of components of $X$.

So we now assume $X$ is a smooth connected curve. Let $\on{Hom}(\pi_1(C^\circ, c), \on{GL}_r(\mathbb{C}))$ be the representation variety parametrizing $r$-dimensional complex representations of $\pi_1(C^\circ, c)$. There is a natural $\on{GL}_r(\mathbb{C})$ action on $\on{Hom}(\pi_1(C^\circ, c), \on{GL}_r(\mathbb{C}))$ induced by conjugation. Let $$f: X\to \on{Hom}(\pi_1(C^\circ, c), \on{GL}_r(\mathbb{C}))$$ be the map sending a point $x$ to $\rho_x|_{\pi_1(C^\circ, c)}$. We wish to show that the image of $f$ lies in a single $\on{GL}_r(\mathbb{C})$-orbit, namely that of $\rho_y|_{\pi_1(C^\circ, c)}$.

Choose a local parameter $t$ of $X$ at $y$, so that $\widehat{\mathscr{O}_{X,y}}\simeq \mathbb{C}[[t]]$. Let $$\rho_\infty: \pi_1(\mathscr{C}^\circ, c)\to GL_r(\mathbb{C}[[t]])$$ be the corresponding representation. By \autoref{lemma:infinitesmal-rigidity}, $\rho_\infty|_{\pi_1(C^\circ, c)}$ is conjugate to a constant representation. It follows that there exists a dense open
subset $U\subset X$ containing $y$ such that $f(U)$ is contained in the
$\on{GL}_r(\mathbb{C})$-orbit of $\rho_y|_{\pi_1(C^\circ, c)}$.

As $\rho_y|_{\pi_1(C^\circ, c)}$ is unitary, hence semisimple, its $\on{GL}_r(\mathbb{C})$-orbit in $\on{Hom}(\pi_1(C^\circ, c), \on{GL}_r(\mathbb{C}))$ is closed \cite[Theorem 30]{sikora2012character}. Hence the closure $\overline{U}$ of $U$ maps into the $\on{GL}_r(\mathbb{C})$-orbit of $\rho_y|_{\pi_1(C^\circ, c)}$. But as $X$ is irreducible, $\overline{U}=X$, completing the proof.
\end{proof}
\begin{theorem}
\source{canonical-representations-surface-groups}
\notready
	\label{theorem:finite-image-semisimple}
\source{canonical-representations-surface-groups}
	Let $$\rho: \pi_1(\Sigma_{g,n})\to \on{GL}_r(\mathbb{C})$$ be a semisimple MCG-finite representation, and suppose $r<\sqrt{g+1}$. Then $\rho$ has finite image.
\end{theorem}
\begin{proof}
	Observe that if $\rho$ is a sum of irreducible representations, each of
	which have finite monodromy, then $\rho$ has finite monodromy as well.
	As a summand of an MCG-finite representation is MCG-finite by
	\autoref{proposition:basic-properties}, it suffices to treat the case that $\rho$ is irreducible.

	By \autoref{corollary:converse-to-MCG-rep}, there exists a dominant
	\'etale map $\mathscr{M}\to \mathscr{M}_{g,n}$, and a local system
	$\mathbb{V}$ with finite determinant on the total space $\mathscr{C}^\circ$ of the associated
	family of $n$-times punctured curves of genus $g$, such that for
	$C^\circ$ a fiber of $\pi^\circ$, $\mathbb{V}|_{{C}^\circ}$ has monodromy representation given by $\rho$. 
		By \autoref{theorem:deformation-to-pvhs}, $\mathbb V$ admits a deformation with constant determinant to a
	local system $\mathbb V_0$ underlying a complex PVHS.
	Since $r < 2\sqrt{g+1}$, it follows from \cite[Theorem
	1.2.12]{LL:geometric-local-systems} that
	the local system $\mathbb V_0|_{{C}^\circ}$ is unitary. Moreover, $\mathbb V_0|_{{C}^\circ}$ is MCG-finite by \autoref{proposition:MCG-finite-family-of-curves}.
	
	It follows from \autoref{proposition:unitary-implies-finite}
	that $\mathbb V_0|_{{C}^\circ}$ has finite monodromy.
	To conclude, it's enough to show that $\mathbb V|_{{C}^\circ} =
	\mathbb V_0|_{{C}^\circ}$.
	But this is immediate from \autoref{lemma:global-rigidity}.
\end{proof}

\subsection{Completion of the non-semisimple case}
\label{subsection:non-semisimple}
\source{canonical-representations-surface-groups}
We come to the final step, where we verify the non-semisimple case of \autoref{theorem:finite-image}.
The idea is to show that if we do have a non-semisimple MCG-finite representation of $\pi_1(\Sigma_{g,n})$ of low rank, we
can produce a certain
 finite cover $\Sigma_{g',n'}$ of $\Sigma_{g,n}$, which violates 
our results toward the Putman-Wieland conjecture.

If $\rho$ is a $G$-representation and $\rho'\subset \rho$ is a subrepresentation, we say that $\rho'$ is \emph{characteristic} if it is stable under all $G$-automorphisms of $\rho$ (for example, the socle of $\rho$, or an isotypic component thereof).
\begin{lemma}
\source{canonical-representations-surface-groups}
\notready
	\label{lemma:semisimplicity}
\source{canonical-representations-surface-groups}
Let $\rho$ be an MCG-finite representation of $\pi_1(\Sigma_{g,n})$, and let $\rho_1\subset\rho$ be a semisimple, characteristic subrepresentation of $\rho$. Let $\rho_2=\rho/\rho_1$, and suppose $\rho_2$ is semisimple as well. If $\rho_1, \rho_2$ have finite image and $\rk\rho<\sqrt{g+1}$, then $\rho$ has finite image as well, i.e.~the extension of $\rho_2$ by $\rho_1$ splits.
\end{lemma}
\begin{proof}
If $g=0$, $\rho$ is the zero representation, so without loss of generality we may assume $g$ is positive.

	Decompose $\rho_2^\vee \otimes \rho_1$ as $\rho_2^\vee \otimes \rho_1 \simeq \oplus_i \sigma_i^{\oplus n_i}$, with the $\sigma_i$ irreducible and pairwise non-isomorphic, and the $n_i$ positive integers.
	Since $\dim \rho_1 + \dim \rho_2 = \dim \rho < \sqrt{g+1}$, 
	we obtain from the AM-GM inequality that $\dim \rho_1 \cdot \dim \rho_2
	< (\sqrt{g+1}/2)^2 = (g+1)/4$. Therefore, for every $i$, $$n_i \cdot \dim
	\sigma_i < (g+1)/4$$
	and hence $n_i, \dim\sigma_i < (g+1)/4$. 
	To put ourselves in the setting of 
	\autoref{theorem:asymptotic-putman-wieland}, we choose an additional
	basepoint $v \in \Sigma_{g,n}$ and let $\Sigma_{g,n+1} := \Sigma_{g,n} \setminus\{v\}$. Let $\Gamma\subset \on{Mod}_{g, n+1}$ be a finite index subgroup stabilizing the conjugacy class of each $\sigma_i$.

	Let $\Sigma_{g',n'} \to \Sigma_{g,n}$ be a finite Galois cover, with Galois group $H$, upon which the local
	systems corresponding to both
	$\rho_1$ and $\rho_2$ trivialize (for example, the cover defined by $\ker(\rho_1\oplus\rho_2)$).
	It suffices to show that $\rho|_{\pi_1(\Sigma_{g',n'})}$ is trivial. As
	$\rho|_{\pi_1(\Sigma_{g',n'})}$ has abelian image, it factors through an
	$H$-equivariant map $$H_1(\Sigma_{g',n'})\to \rho_2^\vee\otimes
	\rho_1.$$ Equivalently, this is the extension class corresponding to
	$\rho$ in $$\on{Ext}^1_{\pi_1(\Sigma_{g', n'})}(\rho_2,
	\rho_1)=H^1(\pi_1(\Sigma_{g', n'}), \Hom(\rho_2,
	\rho_1))=\Hom(\pi_1(\Sigma_{g', n'}), \Hom(\rho_2, \rho_1)).$$
	Projecting onto the $\sigma_i$-isotypic piece of $\rho_2^\vee\otimes
	\rho_1$ yields a $\Gamma$-stable quotient of
	$H_1(\Sigma_{g',n'})^{\sigma_i}$, or equivalently a $\Gamma$-stable
	subspace of $H^1(\Sigma_{g',n'})^{\sigma_i}$, of rank at most $(g+1)/4$.
	Call this subspace $Q_i$.
	But by
	\autoref{theorem:asymptotic-putman-wieland}, we have that any nonzero
	$\Gamma$-stable subspace of $H^1(\Sigma_{g',n'})^{\sigma_i}$ has rank at
	least $2g-2\dim \sigma_i$, which satisfies the inequality
	\begin{align*}
		2g-2\dim \sigma_i >2g-(g+1)/2=\frac{3g}{2}-\frac{1}{2}>(g+1)/4
		> n_i \dim \sigma_i.
	\end{align*}
	Hence $Q_i$ equals zero. As this holds for all $i$, $\rho|_{\pi_1(\Sigma_{g',n'})}$ is trivial as desired.
\end{proof}

\subsection{}
\label{subsection:main-proof}
\source{canonical-representations-surface-groups}
We finally complete the proof of our main theorem, that 
MCG-finite representations $$\rho: \pi_1(\Sigma_{g,n})\to
\on{GL}_r(\mathbb{C})$$ with $r<\sqrt{g+1}$ have finite image.
\begin{proof}[Proof of \autoref{theorem:finite-image}]
The proof is by induction on the length of the socle filtration. The base case is \autoref{theorem:finite-image-semisimple}. Now let $\rho_1$ be the socle of $\rho$.	Because the socle of a representation is characteristic both $\rho_1$
	and $\rho/\rho_1$ are
	 MCG-finite. 
	Therefore, both $\rho_1$ and $\rho/\rho_1$ have finite image by induction.
	We conclude by \autoref{lemma:semisimplicity}.
\end{proof}
